\documentclass{beamer}
\usepackage[utf8]{inputenc}
\usepackage{amsmath}
\usepackage{graphicx}

% Set Beamer Theme
\usetheme{Madrid} % You can change this to another theme

% Title Information
\title[Physics-Constrained Gaussian Splatting]{Research Gaps in Physics-Constrained Gaussian Splatting for Medical Imaging}
\subtitle{Integrating Physical Constraints into 3D Representations}
\author{Your Name}
\institute{Your Institution}
\date{\today}

\begin{document}

% Slide 1: Title Slide
\begin{frame}
    \titlepage
\end{frame}

% Slide 2: Agenda
\begin{frame}{Agenda}
    \tableofcontents
\end{frame}

% Slide 3: Background on Gaussian Splatting
\section{Background on Gaussian Splatting}
\begin{frame}{Background on Gaussian Splatting}
    \begin{itemize}
        \item \textbf{Definition:} Explicit 3D representation using a cloud of Gaussians.
        \item \textbf{Advantages:}
        \begin{itemize}
            \item Real-time rendering.
            \item High efficiency compared to NeRF-based methods.
        \end{itemize}
        \item \textbf{Current Applications:}
        \begin{itemize}
            \item RGB novel view synthesis.
            \item Emerging use in X-ray and CT reconstruction.
        \end{itemize}
    \end{itemize}
\end{frame}

% Slide 4: Current State-of-the-Art in Medical Imaging
\section{Current State-of-the-Art in Medical Imaging}
\begin{frame}{Current State-of-the-Art in Medical Imaging}
    \begin{itemize}
        \item \textbf{Neural Radiance Fields (NeRF):}
        \begin{itemize}
            \item Implicit representation.
            \item High-quality reconstructions.
            \item Limitations: long training/inference times, heavy computation.
        \end{itemize}
        \item \textbf{Gaussian Splatting (3DGS):}
        \begin{itemize}
            \item Fast and efficient in the RGB domain.
            \item Initial adaptations for X-ray (e.g., X-Gaussian).
        \end{itemize}
        \item \textbf{Key Limitation:} Lack of explicit integration of physical imaging models.
    \end{itemize}
\end{frame}

% Slide 5: Identification of the Research Gap
\section{Identification of the Research Gap}
\begin{frame}{Identification of the Research Gap}
    \begin{itemize}
        \item \textbf{The Gap:} Current Gaussian splatting models do not incorporate known physical constraints.
        \item \textbf{Why It Matters:}
        \begin{itemize}
            \item In X-ray imaging: disregard of the Beer–Lambert law.
            \item In MRI: neglect of complex-valued signal formation (T1/T2 relaxations, echo time).
        \end{itemize}
    \end{itemize}
\end{frame}

% Slide 6: Understanding Physical Constraints
\section{Understanding Physical Constraints}
\begin{frame}{Understanding Physical Constraints}
    \begin{itemize}
        \item \textbf{X-ray Attenuation:}
        \[
        I = I_0 \exp\left(-\int_{\text{ray}} \mu(x) \,dx\right)
        \]
        \item \textbf{MRI Signal Model:}
        \begin{itemize}
            \item Signal equations depend on proton density and relaxation times (T1, T2).
        \end{itemize}
    \end{itemize}
\end{frame}

% Slide 7: Exploration Direction 1 – Embedding X-ray Attenuation Laws
\section{Exploration Direction 1}
\begin{frame}{Embedding X-ray Attenuation Laws}
    \begin{itemize}
        \item \textbf{Redesigning the Gaussian Kernel:}
        \begin{itemize}
            \item Associate each Gaussian with an attenuation (density) parameter.
        \end{itemize}
        \item \textbf{Physics-Based Loss Function:}
        \[
        \mathcal{L}_{\text{physics}} = \| I_{\text{measured}} - I_0 \exp(-\sum_{i} \mu_i \, w_i) \|_1
        \]
    \end{itemize}
\end{frame}

% Slide 8: Exploration Direction 2 – Incorporating MRI Signal Models
\section{Exploration Direction 2}
\begin{frame}{Incorporating MRI Signal Models}
    \begin{itemize}
        \item \textbf{Adapting Gaussian Splatting for MRI:}
        \begin{itemize}
            \item Extend representation to handle complex signals.
        \end{itemize}
        \item \textbf{Differentiable MRI Forward Model:}
        \begin{itemize}
            \item Simulate MRI acquisition (Fourier sampling, signal decay).
        \end{itemize}
    \end{itemize}
\end{frame}

% Slide 9: Benefits of a Unified, Physics-Constrained Framework
\section{Benefits of a Unified Framework}
\begin{frame}{Benefits of a Unified, Physics-Constrained Framework}
    \begin{itemize}
        \item \textbf{Accuracy & Robustness:} Reduces artifacts and noise.
        \item \textbf{Interpretability:} Parameters correspond to actual physical quantities.
        \item \textbf{Clinical Impact:} Potential for lower radiation dose and better diagnostic quality.
    \end{itemize}
\end{frame}

% Slide 10: Roadmap for Implementation
\section{Roadmap for Implementation}
\begin{frame}{Roadmap for Implementation}
    \begin{itemize}
        \item \textbf{Step 1: Theoretical Derivation}
        \item \textbf{Step 2: Kernel Redesign}
        \item \textbf{Step 3: Loss Function Engineering}
        \item \textbf{Step 4: Initialization & Regularization}
        \item \textbf{Step 5: Validation & Testing}
        \item \textbf{Step 6: Optimization & Scalability}
    \end{itemize}
\end{frame}

% Slide 11: Challenges and Future Directions
\section{Challenges and Future Directions}
\begin{frame}{Challenges and Future Directions}
    \begin{itemize}
        \item \textbf{Computational Complexity:} Balancing real-time performance.
        \item \textbf{Validation:} Need for extensive clinical testing.
        \item \textbf{Extension to Multimodal Data:} Integrating PET/MRI.
    \end{itemize}
\end{frame}

% Slide 12: Summary and Conclusion
\section{Summary and Conclusion}
\begin{frame}{Summary and Conclusion}
    \begin{itemize}
        \item \textbf{Recap:} Gaussian splatting lacks physics-based constraints.
        \item \textbf{Proposed Exploration:} Embedding X-ray and MRI models.
        \item \textbf{Future Impact:} Improved reconstruction quality, clinical trust.
    \end{itemize}
\end{frame}

\end{document}